\documentclass{article}
\usepackage{mathtools} 
\usepackage{fontspec}
\usepackage[UTF8]{ctex}
\usepackage{amsthm}
\usepackage{mdframed}
\usepackage{xcolor}
\usepackage{amssymb}
\usepackage{amsmath}


% 定义新的带灰色背景的说明环境 zremark
\newmdtheoremenv[
  backgroundcolor=gray!10,
  % 边框与背景一致，边框线会消失
  linecolor=gray!10
]{zremark}{说明}

% 通用矩阵命令: \flexmatrix{矩阵名}{元素符号}{行数}{列数}
\newcommand{\flexmatrix}[4]{
  \[
  #1 = \begin{pmatrix}
    #2_{11}     & #2_{12}     & \cdots & #2_{1#4}   \\
    #2_{21}     & #2_{22}     & \cdots & #2_{2#4}   \\
    \vdots      & \vdots      & \ddots & \vdots     \\
    #2_{#31}    & #2_{#32}    & \cdots & #2_{#3#4}
  \end{pmatrix}
  \]
}

% 简化版命令（默认矩阵名为A，元素符号为a）: \quickmatrix{行数}{列数}
\newcommand{\quickmatrix}[2]{\flexmatrix{A}{a}{#1}{#2}}

\begin{document}
\title{习题 13.2}
\author{张志聪}
\maketitle

\section*{1}
\begin{itemize}
  \item  (2)
        (a-i)
        任意$x \in [0, 1]$
        \begin{align*}
          \lim\limits_{n \to +\infty} x - \frac{x^n}{n}
           & = x
        \end{align*}
        令$f(x) = x$，于是$f(x)$是函数列的极限函数。
        \begin{align*}
          \sup\limits_{x \in [0, 1]} |f_n(x) - f(x)|
           & = \sup\limits_{x \in [0, 1]} \left|\frac{x^n}{n} \right| \\
           & = \frac{1}{n}
        \end{align*}
        于是
        \begin{align*}
          \lim\limits_{n \to +\infty} \sup\limits_{x \in [0, 1]} |f_n(x) - f(x)|
           & = \lim\limits_{n \to +\infty} \frac{1}{n} \\
           & = 0
        \end{align*}
        故$\{f_n(x)\}$一致收敛于$f(x)$。

        (a-ii)
        \begin{align*}
          f_n'(x) = 1 - x^{n - 1}
        \end{align*}
        任意$x \in [0, 1)$
        \begin{align*}
          \lim\limits_{n \to +\infty} f_n'(x)
           & = \lim\limits_{n \to +\infty} 1 - x^{n - 1} \\
           & = 1
        \end{align*}
        $x = 1$时
        \begin{align*}
          \lim\limits_{n \to +\infty} f_n'(x)
           & = \lim\limits_{n \to +\infty} 1 - x^{n - 1} \\
           & = 0
        \end{align*}
        令
        \begin{equation*}
          f'(x) = \begin{cases}
            1, & x \in [0, 1) \\
            0, & x = 1
          \end{cases}
        \end{equation*}
        于是$f'(x)$是函数列$\{f_n'(x)\}$的极限函数。

        $x \in [0, 1)$时
        \begin{align*}
          \lim\limits_{x \to 1^-} |f_n'(x) - f'(x)|
           & = \lim\limits_{x \to 1} x^{n - 1} \\
           & = 1
        \end{align*}
        由保号性可知
        \begin{align*}
          \sup\limits_{x \to [0, 1)} |f_n'(x) - f'(x)| > \frac{1}{2}
        \end{align*}
        另外，$x = 1$时，不影响该性质。

        所以
        \begin{align*}
          \lim\limits_{n \to +\infty} \sup\limits_{x \in [0, 1]} |f_n'(x) - f'(x)|
           & \neq 0
        \end{align*}
        故$\{f_n'(x)\}$不一致收敛于$f'(x)$。

        \textcolor{red}{另证：}
        因为$f_n'(x)$连续，而$f'(x)$不连续，由定理13.9的逆否命题可知，
        $\{f_n'(x)\}$在$[0, 1]$上不一致。

        (b) 满足13.9、13.10的条件和结论。不满足13.11的条件，接下来证明其结论。

        \begin{align*}
          \frac{d}{dx} \left(\lim\limits_{n \to +\infty} f_n(x) \right)
           & = \frac{d}{dx} f(x) \\
           & = \frac{d}{dx} x    \\
           & = 1
        \end{align*}
        又有之前的讨论知
        \begin{align*}
          \lim\limits_{n \to +\infty} \frac{d}{dx} f_n(x)
           & = f'(x)
        \end{align*}
        所以
        \begin{align*}
          f'(x) \neq 1
        \end{align*}
        故13.11的结论不成立。

  \item (3)

        (a-i)
        \begin{align*}
          \lim\limits_{n \to +\infty} f_n(x)
           & = \lim\limits_{n \to +\infty} \frac{nx}{e^{nx^2}}
        \end{align*}
        $x \neq 0$，由洛必达法则可知（注：$n$看做变量，$x$看做变量）
        \begin{align*}
          \lim\limits_{n \to +\infty} \frac{nx}{e^{nx^2}} = 0
        \end{align*}
        $x = 0$时
        \begin{align*}
          \lim\limits_{n \to +\infty} \frac{nx}{e^{nx^2}}
           & = \lim\limits_{n \to +\infty} 0 \\
           & = 0
        \end{align*}
        令$f(x) = 0$为函数列的极限函数。
        \begin{align*}
          \sup\limits_{x \in [0, 1]} |f_n(x) - f(x)|
           & = \sup\limits_{x \in [0, 1]} \frac{nx}{e^{nx^2}}
        \end{align*}
        \begin{align*}
          f_n'(x) & = \left(\frac{nx}{e^{nx^2}} \right)                  \\
                  & = \frac{n e^{nx^2} - nx e^{nx^2}(2nx)}{(e^{nx^2})^2} \\
                  & = \frac{(n - 2nx^2) e^{nx^2}}{(e^{nx^2})^2}          \\
                  & = \frac{n - 2nx^2}{e^{nx^2}}
        \end{align*}
        令
        \begin{align*}
          f_n'(x)                    & = 0                   \\
          \frac{n - 2nx^2}{e^{nx^2}} & = 0                   \\
          n - 2nx^2                  & = 0                   \\
          x                          & = \sqrt{\frac{1}{2n}}
        \end{align*}
        按照$[0, \sqrt{\frac{1}{2n}}], f_n'(x) \geq 0; [\sqrt{\frac{1}{2n}}, 1] \leq 0$，
        所以$x = \sqrt{\frac{1}{2n}}$是极大值点，所以
        \begin{align*}
          \sup\limits_{x \in [0, 1]} \frac{nx}{e^{nx^2}}
           & = \frac{\sqrt{n}}{\sqrt{2 e}}
        \end{align*}
        于是
        \begin{align*}
          \lim\limits_{n \to +\infty} \sup\limits_{x \in [0, 1]} \frac{nx}{e^{nx^2}}
           & = \lim\limits_{n \to +\infty} \frac{\sqrt{n}}{\sqrt{2 e}} \\
           & = + \infty
        \end{align*}
        故$\{f_n(x)\}$不一致收敛。

        (a-ii)
        $x \in (0, 1]$时
        \begin{align*}
          \lim\limits_{n \to +\infty} f_n'(x)
           & = \lim\limits_{n \to +\infty} \frac{n - 2nx^2}{e^{nx^2}}   \\
           & = \lim\limits_{n \to +\infty} \frac{n(1 - 2x^2)}{e^{nx^2}} \\
           & = 0
        \end{align*}
        $x = 0$时
        \begin{align*}
          \lim\limits_{n \to +\infty} f_n'(x)
           & = \lim\limits_{n \to +\infty} n \\
           & = +\infty
        \end{align*}
        令
        \begin{equation*}
          f'(x) = \begin{cases}
            0,       & x \in (0, 1] \\
            +\infty, & x = 0
          \end{cases}
        \end{equation*}
        所以$\{f_n'(x)\}$在$[0, 1]$上不收敛，
        从而也不一致收敛。

        (b-i) 不满足13.9的条件，$f(x) = 0$，故结论成立。

        (b-ii) 不满足13.10的条件。
        \begin{align*}
          \int_{0}^{1} f(x) dx
           & = \int_{0}^{1} 0 dx \\
           & = 0
        \end{align*}
        又有
        \begin{align*}
          \lim\limits_{n \to +\infty} \int_{0}^{1} f_n(x) dx
           & = \lim\limits_{n \to +\infty} \int_{0}^{1} \frac{nx}{e^{nx^2}} dx  \\
           & = \lim\limits_{n \to +\infty} - \frac{1}{2}e^{-nx^2}\bigg|_{0}^{1} \\
           & = \lim\limits_{n \to +\infty} \frac{1}{2} - \frac{1}{2}e^{-n}      \\
           & = \frac{1}{2}
        \end{align*}
        所以
        \begin{align*}
          \int_{0}^{1} f(x) dx \neq \lim\limits_{n \to +\infty} \int_{0}^{1} f_n(x) dx
        \end{align*}
        故结论不成立。

        (b-iii)不满足13.11的条件，接下来证明其结论。

        \begin{align*}
          \frac{d}{dx} \lim\limits_{n \to \infty}\left(f_n(x)\right)
           & = \frac{d}{dx} 0 \\
           & = 0
        \end{align*}
        由之前的证明可得
        \begin{align*}
          0 \neq \lim\limits_{n \to +\infty} \frac{d}{dx} f_n(x) = f'(x)
        \end{align*}
        故13.11的结论不成立。
\end{itemize}

\section*{2}

由于$f_n(x_0) \to f(x_0) \, (n \to \infty), f'_n \rightrightarrows g \, (n \to \infty)$，
因此对任意$\epsilon > 0$，存在正数$N$，当$n > N$时，任意$x \in [a, b]$
\begin{align*}
  |f_n(x_0) - f(x_0)| < \epsilon
\end{align*}
和
\begin{align*}
  |f_n'(x) - g(x)| < \epsilon
\end{align*}
由定理13.11的证明过程可知，
\begin{align*}
  |f_n(x) - f(x)|
   & = \left|f_n(x_0) + \int_{x_0}^{x} f_n'(t) dt - f(x_0) - \int_{x_0}^{x} g(t) dt \right|                   \\
   & \leq \left|f_n(x_0) - f(x_0) \right| + \left| \int_{x_0}^{x} f_n'(t) dt - \int_{x_0}^{x} g(t) dt \right| \\
   & = \left|f_n(x_0) - f(x_0) \right| + \left| \int_{x_0}^{x} f_n'(t) - g(t) dt \right|                      \\
   & \leq \epsilon + \epsilon (b - a)                                                                         \\
   & = (b - a + 1)\epsilon
\end{align*}
由$b - a + 1$是定值和$\epsilon$的任意性可知，
命题得证。

\section*{3}
设$\sum u_n(x)$的部分和函数列为$\{S_n(x)\}$， $\sum u_n(x)$的和函数为$S(x)$；
\begin{itemize}
  \item (1) 定理13.2（连续性）

        \begin{itemize}
          \item 方法一

                因为$\sum u_n(x)$的每一项都连续，所以
                $\{S_n(x)\}$的每一项也都是连续的。
                即
                \begin{align*}
                  \lim\limits_{n \to +\infty} S_n(x) = S(x)
                \end{align*}
                由定理13.9可知，$S(x)$是连续的。

          \item 方法二

                利用极限交换定理（13-2-comment.tex中有证明），仿照定理13.9的证明。

                设$x_0$为$[a, b]$上任一点。由于$\lim\limits_{x \to x_0} u_n(x) = u_n(x_0)$，
                于是由极限交换定理知$\lim\limits_{n \to +\infty} \sum u_n(x)$，且
                \begin{align*}
                  \lim\limits_{x \to x_0} S(x)
                   & = \lim\limits_{x \to x_0} \sum u_n(x) \\
                   & = \sum \lim\limits_{x \to x_0} u_n(x) \\
                   & = \sum u_n(x_0)                       \\
                   & = S(x_0)
                \end{align*}
                因此$S(x)$在$x_0$处连续。

        \end{itemize}

  \item (2) 定理13.4（逐项求导）

        仿照定理13.11的证明。

        设$S_n(x_0) \to A \, (n \to \infty)$，$\sum u_n'(x) \rightrightarrows g(x) \, (n \to \infty), x \in [a, b]$。
        我们要证明函数项级数$\sum u_n(x)$在$[a,b]$上收敛，且其极限函数的导数存在且等于g.

        由定理条件可知，$u_n(x)$是连续函数，对任意$x \in [a, b]$，总有
        \begin{align*}
          u_n(x)
           & = u_n(x_0) + \int_{x_0}^{x} u_n'(t) dt
        \end{align*}
        等式两边求和
        \begin{align*}
          \sum u_n(x)
           & = \sum \left( u_n(x_0) + \int_{x_0}^{x} u_n'(t) dt \right) \\
           & = \sum u_n(x_0) + \sum \int_{x_0}^{x} u_n'(t) dt           \\
           & = A + \int_{x_0}^{x} \sum u_n'(t) dt                       \\
           & = A + \int_{x_0}^{x} g(t) dt                               
        \end{align*}
        由$g$的连续性以及微积分学基本定理推得
        \begin{align*}
          \frac{d}{dx} \left(\sum u_n(x)\right) = g(x) = \sum (\frac{d}{dx} u_n(x))
        \end{align*}




\end{itemize}



\end{document}